\newpage
\section{误差分析}
\subsection{状态评估模型误差}
熵权法权重分配存在主观性偏差，不同年份数据可能导致权重波动。通过Bootstrap重采样（1000次）估计权重标准差，发现主要气体（H2、C2H2、C2H4）权重稳定性最高（CV<5%）。健康/亚健康/故障三级分类的混淆矩阵显示，亚健康样本有时被误判为健康（5.7%），这是正常现象。
\subsection{故障诊断模型误差}
CNN-LSTM融合模型的误差主要来源于：
\begin{enumerate}
\item[(1)] \textbf{多源数据时间对齐误差}:局部放电（微秒级）与油色谱（小时级）采样频率差异大，需要插值对齐，引入额外噪声。
\item[(2)] \textbf{特征提取局限性}:时域/频域特征对早期微弱故障不敏感，部分潜伏期故障的DGA特征尚未显现。
\item[(3)] \textbf{类别不平衡}:故障样本远少于健康样本，采用SMOTE过采样后少数类F1从0.62提升至0.81。
\end{enumerate}
灵敏度分析表明，油色谱数据缺失率<15%时，诊断准确率下降不超过3个百分点。
\subsection{RUL预测误差}
Gamma过程建模的RUL预测存在置信区间较宽的问题。95%置信区间宽度约为点估计值的30%。主要误差来源是退化轨迹的个体差异性。
